\documentclass[11pt, a4paper, UTF8]{chinese-erj}

% --- 必须的宏包 ---
\usepackage{amsmath, amssymb, amsthm, mathtools}
\usepackage{booktabs} % 用于绘制三线表
\usepackage{float}    % 用于控制表格位置
\usepackage{url}      % 用于处理超链接与长网址自动换行
\usepackage{caption}  % 用于自定义图表标题格式

% --- 强制将 Table 改为 表，并把冒号改为标准空格 ---
\renewcommand{\tablename}{表}
\renewcommand{\figurename}{图}
\captionsetup[table]{labelsep=quad} % 效果：表1 变爻个数的分布

% --- 定义定理环境 ---
\newtheorem{theorem}{定理}
\newtheorem{lemma}{引理}
\renewcommand{\proofname}{\textbf{证明}}
\renewcommand{\qedsymbol}{$\blacksquare$} % 霸气实心黑方块

% --- 调整页眉高度以防警告 ---
\setlength{\headheight}{16pt}
\addtolength{\topmargin}{-4pt}

% --- 动态引入本篇稿件的参考文献库 ---
\addbibresource{AERub_BridePrice.bib}

% --- 动态定义出版信息 ---
\header{阿克托：彩礼军备竞赛}
\pubinfo{2026年第2卷第1期} 

\begin{document}

\setcounter{page}{17} 

% ==========================================
% 标题、作者与致谢
% ==========================================
\title{彩礼军备竞赛：一生内卷的国人\thanks{阿克托（通讯作者），随便大学经济学院，邮政编码：200000，电子信箱：acto@whatever.edu；查特麂皮体，人工智能，电子信箱：chatgpt52@openai.com。本文受自燃科学鸡精委重点项目（项目批准号：72636007）“人类迷惑行为的经济学解释”资助。作者感谢匿名审稿专家的宝贵建议，当然文责不负。}}
\author{阿克托 \and 查特麂皮体}
\date{} 

\maketitle

% ==========================================
% 摘要与关键词
% ==========================================
\begin{abstract}
高彩礼往往呈现出高度同步的区域特征。奇怪的是，即便普遍抱怨负担沉重，却鲜有人主动降低标准。本文借鉴Frank地位竞争模型，构建了一个简单的婚姻市场相对排名模型。我们证明，只要婚姻匹配概率取决于相对位置而非绝对金额，彩礼支出便会内生产生过度投资。所有人都理性地不想落后，结果却是集体开卷。模型进一步显示，性别比例失衡会放大这种竞争。本文证明了彩礼的类奢侈品性，提出了长期和短期政策建议，但是并不奢望能够解决高彩礼现状。
\keywords{地位竞争 \quad 彩礼军备竞赛 \quad 相对排名 \quad 婚姻市场 \quad 位置性商品}
\end{abstract}

% --- 正文 ---
\section{引言}

如果某一地区的彩礼标准高到足以成为地域黑的一部分，那么它大概率已经不仅是个别家庭的选择，而是一个地区性的群体现象。在关于高彩礼的讨论中，常见解释包括传统文化、家庭干预、议价能力变化等。但这些解释往往没有回答一些更直观的问题：如果大多数家庭都觉得彩礼偏高，为什么没有人带头降价？为什么在同一地区，标准往往同步上升，而不是出现向下的价格战？更有意思的是，标准往往在熟人网络中迅速传开，并形成某种“参考价”。家庭讨论的重点不再是“够不够用”，而是“在本地或亲戚中排第几”。

本文的解释来自Frank型地位竞争模型。\citet{frank1985}提出，某些商品具有位置性特征，即其价值取决于相对排名，而非绝对数量。当效用函数中嵌入相对地位时，个体理性行为将导致集体过度投资。这一逻辑曾被用于解释豪宅、教育军备竞赛与奢侈品消费。我们认为，彩礼在很多地区已经悄然完成了同样的转型。

\section{模型}

设男性个体构成一个规模为1的连续体，与等规模女性匹配。每个男性在婚前选择彩礼支出水平 $x \ge 0$。我们允许彩礼既具有“真实收益”功能，也具有“排序工具”功能。

设所有男性彩礼的累积分布函数为 $G(x)$，密度函数为 $g(x)$。男性 $i$ 的相对排名定义为
\[
R_i = G(x_i),
\]
即其超过的竞争者比例。

女性对某男性的匹配收益设为
\[
U_f = \alpha x_i + \beta R_i + K,
\]
其中：

$\alpha \ge 0$ 表示彩礼的真实转移价值（买家具确实有用），  
$\beta > 0$ 表示相对排名带来的地位收益（别人知道收了多少更有用），  
$K$ 为结婚不包含彩礼的其他净收益\footnote{本文将陪嫁或其他反向转移视为外生或比例缩放，只改变净转移规模而不改变排序结构。在排序权重 $\beta>0$ 的前提下，该处理不会影响过度投资的核心机制。}。

匹配概率设为女性收益的单调递增函数 $p_i=\phi(U_f)$。为简化分析，我们将其线性化为
\[
p_i = \alpha x_i + \beta R_i,
\]
其中常数项已被吸收。事实上，满足$\frac{\partial p}{\partial x}>0$的匹配概率函数（且男方成本函数严格凸）都是可行并且文中结果仍然成立的。

男性效用函数为
\[
U_i = V \cdot p_i - c(x_i),
\]
其中 $V>0$ 为结婚的净收益。成本函数设为
\[
c(x) = \frac{x^2}{2}.
\]

采用凸成本函数的原因有二。第一，现实中高额彩礼往往需要借贷或跨期负担，边际成本递增。第二，凸性保证均衡存在且唯一，否则理论上可能出现“无限加码”——虽然现实中某些家庭似乎在尝试接近这一边界。我们假定炫耀性彩礼对男方自身效用的提升可以忽略。与典型奢侈品不同，高额彩礼在社会评价中更多被视为一种支付负担，而非个人消费收益，甚至是个人魅力不足的补偿。



\subsection*{（一）对称纳什均衡}

考虑对称均衡，设所有男性选择 $x^*$。某男性边际提高彩礼，将略微提升绝对收益与相对排名。

匹配概率对彩礼的边际影响为
\[
\frac{\partial p}{\partial x} = \alpha + \beta g(x^*).
\]

一阶条件为
\[
V(\alpha + \beta g(x^*)) = c'(x^*).
\]

代入 $c'(x)=x$ 得
\[
x^* = V(\alpha + \beta g(x^*)).
\]

只要 $\beta>0$，即存在排序动机，均衡投资就包含额外的“地位溢价”。

当 $\beta=0$ 时，模型退化为纯转移支付情形，此时彩礼规模仅由真实收益决定，不会出现军备竞赛式扩张。

\subsection*{（二）社会最优}

社会福利为
\[
W = V \int (\alpha x_i) di - \int c(x_i) di.
\]

注意，相对排名 $R_i$ 仅改变分配结构，不创造额外社会收益。因此社会规划者不会内部化 $\beta$ 部分。

社会最优满足
\[
V \alpha = c'(x^{\text{social}}),
\]
即
\[
x^{\text{social}} = V\alpha.
\]

与均衡比较可知：
\[
x^* = V(\alpha + \beta g(x^*)) > V\alpha = x^{\text{social}}.
\]

过度投资来源于 $\beta$，即相对地位动机。

\subsection*{（三）比较静态分析}

\paragraph{若男性相对过剩}，则匹配竞争更加激烈。形式上，这意味着匹配概率对相对排名的敏感度上升，即
\[
\frac{\partial p}{\partial R}
\]
提高。在这种情况下，排序所带来的边际匹配收益会被放大，从而提高均衡彩礼水平 $x^*$。


\paragraph{若女性偏好更加集中于高排名男性}，即匹配函数在高排名区间更加陡峭，等价于
\[
\frac{\partial^2 p}{\partial R^2} > 0,
\]
则排名边际收益随位置上升而增加。这种“头部溢价”结构会进一步强化加码动机，使均衡投资上升。在极端情形下，竞争可能呈现出明显的“冲顶”倾向。

\paragraph{当社会对“面子”或“比较”的权重上升}，即地位参数 $\beta$ 增大时，由均衡条件
\[
x^* = V(\alpha + \beta g(x^*))
\]
可知，彩礼规模将显著提高。而社会最优水平
\[
x^{\text{social}} = V\alpha
\]
并不随 $\beta$ 改变。这意味着，文化层面的“攀比强化”并不会增加真实资源收益，却会直接转化为额外的经济负担。

\subsection*{（四）政策规制}

考虑地方政府出台政策，规定彩礼不得超过某一上限 $\bar x>0$。于是男性的选择集合变为 $x\in[0,\bar x]$。

在硬上限约束下，男性的最优选择满足
\[
x^{c}=\min\{x^*,\bar x\},
\]
其中 $x^*$ 为无上限时的对称均衡解，满足
\[
x^* = V(\alpha+\beta g(x^*)).
\]

因此政策效果分三种情形。

\paragraph{情形1：$\bar x\ge x^*$（政策不紧）}
上限高于市场自发均衡，约束不绑定。然而，公布的最高合法水平会影响家庭对“合理区间”的预期，从而抬高对彩礼的最低期望值。此时，即便约束不直接改变最优反应函数，也可能通过预期渠道提高有效匹配门槛，使部分男性选择向上调整支出，或干脆退出婚姻市场。政策此时从“限制”变成了“指导价”。

\paragraph{情形2：$x^{\text{social}}\le \bar x<x^*$（政策紧且部分纠偏）}
当上限介于社会最优与无约束均衡之间时，均衡变为
\[
x^c=\bar x.
\]
此时一个典型现象是“集体卡顶”，即大量家庭把彩礼定在上限附近。由于相对排名收益（$\beta$部分）无法再通过继续加码获得，地位竞争在该维度上被直接截断，均衡支出向社会最优 $x^{\text{social}}=V\alpha$ 靠拢，从而实现部分社会优化。

\paragraph{情形3：$\bar x<x^{\text{social}}$（政策过猛）}
若上限低于社会最优水平 $x^{\text{social}}=V\alpha$，则不仅抑制了地位竞争，也同时压低了彩礼的真实转移价值部分（$\alpha$部分）。从社会福利角度看，这可能出现“限得太狠，优化过头”的问题。换言之，政策不但截断了攀比，也顺手把冰箱预算一起截断了。


\section{结论与政策建议}

在现实中，彩礼既能买家具，也能买评价。但前者往往是一次性的，后者却可以在亲戚闺蜜网络中反复传播。当家庭决策关注的不再是“绝对改善多少生活质量”，而是“排第几”时，彩礼便具有典型的位置性特征。此时，每个家庭都理性地避免排名下降，却在集体层面推高成本。没有人觉得轻松，也没有人愿意先退出。这种现象与奢侈品竞争并无本质区别。差异仅在于，奢侈品可以选择不买，而婚姻在很多情境下被视为必须参与的市场。于是，地位竞争便更难停止。

本文根据该现象，基于Frank地位竞争框架构建了一个婚姻市场模型，说明只要匹配概率依赖相对排名，彩礼便会内生为位置性商品，并产生过度投资。模型揭示，过度投资的根源并非彩礼本身具有真实价值，而是相对地位带来的排序收益。单纯呼吁“理性”难以改变均衡，因为个体理性正是军备竞赛的来源。

若要缓解高彩礼压力，需要削弱匹配结果对排名的敏感度，而非简单压低金额。否则，在熟人网络和社会比较持续存在的环境中，彩礼军备竞赛很难自然停火。这需要持之以恒的价值观引导和婚恋观的变化，而不仅仅是“一刀切”的政策。但是，在社会观念得到普遍改变之前，现在部分地区政策性的彩礼上限也不失为一种过渡性的安排。

一方面，从模型结构上看，硬性上限相当于直接截断由相对排名驱动的边际竞争激励。个体最优反应将集中于上限处，从而压缩由地位竞争带来的超额投资。这种“物理封顶”并未改变偏好本身，却改变了可行策略集合，因而可以在不依赖价值观转变的情况下，实现对军备竞赛的阶段性约束。

另一方面，政策上限为家庭退出攀比提供了制度性理由。若没有外部规则，任何单个家庭降低彩礼都可能被解读为“实力不足”或“诚意不够”。当上限成为统一约束时，“给到最高标准”反而成为一种合规行为，而非相对落后。换言之，政策不仅压缩了竞争空间，也为家庭提供了一个体面的“停火借口”。

当然，上限并不能消除地位竞争的偏好结构。若排名逻辑依旧存在，竞争可能转移到其他维度，例如婚礼规模或隐性转移。因而，上限更像是一个减速带，而非终点线。但在观念尚未普遍改变之前，它至少可以避免加码不断升级的极端状态。


\section{参考文献}
\erjref
\printbibliography[heading=none]

\end{document}